\documentclass[10pt,a4paper]{amsart}
\usepackage{amsmath,amssymb,amsthm,graphicx}
% Retained preamble for the Stacks Project chapter snapshot in this directory.
% The chapter's own preamble.tex is a copy of the shared Stacks Project preamble; the block
% below is sliced MECHANICALLY, line for line, from it. Nothing was added, removed or
% redefined; the wrapper class, the xr-hyper/hyperref wiring and the externaldocument list
% are omitted because the assembler supplies the document class.
%
% lines taken from preamble.tex: 3, 5, 9, 11, 12, 156, 157, 158, ... (51 lines total)
%
% package coverage restored from the chapter's own preamble (2026-10-13): xy, xy-2cell

\usepackage[all]{xy}
\xyoption{2cell}
\UseAllTwocells

\newenvironment{reference}{\comment}{\endcomment}
\newenvironment{slogan}{\comment}{\endcomment}
\newenvironment{history}{\comment}{\endcomment}
\theoremstyle{plain}
\newtheorem{theorem}[subsection]{Theorem}
\newtheorem{proposition}[subsection]{Proposition}
\newtheorem{lemma}[subsection]{Lemma}
\theoremstyle{definition}
\newtheorem{definition}[subsection]{Definition}
\newtheorem{example}[subsection]{Example}
\newtheorem{exercise}[subsection]{Exercise}
\newtheorem{situation}[subsection]{Situation}
\theoremstyle{remark}
\newtheorem{remark}[subsection]{Remark}
\newtheorem{remarks}[subsection]{Remarks}
\numberwithin{equation}{subsection}
\def\lim{\mathop{\mathrm{lim}}\nolimits}
\def\colim{\mathop{\mathrm{colim}}\nolimits}
\def\Spec{\mathop{\mathrm{Spec}}}
\def\Hom{\mathop{\mathrm{Hom}}\nolimits}
\def\Ext{\mathop{\mathrm{Ext}}\nolimits}
\def\SheafHom{\mathop{\mathcal{H}\!\mathit{om}}\nolimits}
\def\SheafExt{\mathop{\mathcal{E}\!\mathit{xt}}\nolimits}
\def\Sch{\mathit{Sch}}
\def\Mor{\mathop{\mathrm{Mor}}\nolimits}
\def\Ob{\mathop{\mathrm{Ob}}\nolimits}
\def\Sh{\mathop{\mathit{Sh}}\nolimits}
\def\NL{\mathop{N\!L}\nolimits}
\def\CH{\mathop{\mathrm{CH}}\nolimits}
\def\proetale{{pro\text{-}\acute{e}tale}}
\def\etale{{\acute{e}tale}}
\def\QCoh{\mathit{QCoh}}
\def\Ker{\mathop{\mathrm{Ker}}}
\def\Im{\mathop{\mathrm{Im}}}
\def\Coker{\mathop{\mathrm{Coker}}}
\def\Coim{\mathop{\mathrm{Coim}}}
\DeclareMathSymbol{\boxtimes}{\mathbin}{AMSa}{"02}
\def\QCohstack{\mathcal{QC}\!\mathit{oh}}
\def\Cohstack{\mathcal{C}\!\mathit{oh}}
\def\Spacesstack{\mathcal{S}\!\mathit{paces}}
\def\Quotfunctor{\mathrm{Quot}}
\def\Hilbfunctor{\mathrm{Hilb}}
\def\Curvesstack{\mathcal{C}\!\mathit{urves}}
\def\Polarizedstack{\mathcal{P}\!\mathit{olarized}}
\def\Complexesstack{\mathcal{C}\!\mathit{omplexes}}
\def\Pic{\mathop{\mathrm{Pic}}\nolimits}
\def\Picardstack{\mathcal{P}\!\mathit{ic}}
\def\Picardfunctor{\mathrm{Pic}}
\def\Deformationcategory{\mathcal{D}\!\mathit{ef}}
\usepackage{cleveref}
\crefname{originallemma}{original lemma}{original lemmas}
\Crefname{originallemma}{Original lemma}{Original lemmas}
\makeatletter
\newlabel{lemma-free-graded-module-good}{{0FSG}{source}{}{sourcepointer-1}{}}
\newlabel{lemma-free-graded-module-good@cref}{{[originallemma][0][]0FSG}{[1][0][]source}{}{}{}}
\makeatother
\crefname{originallemma}{original lemma}{original lemmas}
\Crefname{originallemma}{Original lemma}{Original lemmas}
\makeatletter
\newlabel{lemma-good-admissible-ses}{{0FSC}{source}{}{sourcepointer-2}{}}
\newlabel{lemma-good-admissible-ses@cref}{{[originallemma][0][]0FSC}{[1][0][]source}{}{}{}}
\makeatother
\crefname{originallemma}{original lemma}{original lemmas}
\Crefname{originallemma}{Original lemma}{Original lemmas}
\makeatletter
\newlabel{lemma-good-direct-sum}{{0FSD}{source}{}{sourcepointer-3}{}}
\newlabel{lemma-good-direct-sum@cref}{{[originallemma][0][]0FSD}{[1][0][]source}{}{}{}}
\makeatother
\crefname{originallemma}{original lemma}{original lemmas}
\Crefname{originallemma}{Original lemma}{Original lemmas}
\makeatletter
\newlabel{remark-sheaf-graded-sets}{{0FSF}{source}{}{sourcepointer-4}{}}
\newlabel{remark-sheaf-graded-sets@cref}{{[originallemma][0][]0FSF}{[1][0][]source}{}{}{}}
\makeatother
\title{Stacks, Differential Graded Sheaves: lemma-resolve}
\author{The Stacks Project Authors}
\date{stacks-project:sdga.tex:local-snapshot-20261012; GNU FDL 1.2 or later}
\begin{document}
\maketitle
\noindent\textit{Editorial note.} Statement span is lines 3185-3197 of sources/stacks-sdga/pinned-original-source.tex; proof span is lines 3199-3261, the complete author proof (61 body lines).\par\medskip
\noindent\textit{Editorial note.} Navigation stubs: 4 cited-but-not-included labels are supplied as non-load-bearing pointers with their pinned tags where the snapshot tag index has them.\par\medskip
\noindent\textit{Editorial note.} extra\_packages is empty on purpose: the retained preamble already loads xy with options and a second declaration clashes in the assembler wrapper.\par\medskip
\noindent\textit{Editorial note.} Text is retained verbatim, including the author's own typos and placeholders.\par\medskip
\noindent\textit{Editorial note.} Queue index 6 of queue\_sdga.json (pb=2700, stmt 3185-3197, proof 3199-3261).\par\medskip
\noindent\textit{Editorial note.} Section 10.2 checked: the region carries a second complete proof environment (the author's 'Second proof', lines 3263-3442, building the resolution as a total complex). Only the FIRST proof (lines 3199-3261) is retained inside the card span, per the gate's one-proof-per-card rule and this lane's precedent; the second proof is an independent alternative and is not counted as a separate card (same chapter + same label).\par\medskip
\noindent\textit{Editorial note.} Cross-chapter same-label warning (NOT an occupancy conflict, Chief ruling 2026-10-13 evening): label 'lemma-resolve' also exists in stacks-dga — a different proposition in a different chapter file.\par\medskip
\noindent\textit{Allowed standard prerequisites (editorial).}\begin{itemize}
\item Explicit prerequisite boundary. The cited inputs are prerequisites; the argument in the retained span is not replaced by any cited result.
\end{itemize}\medskip
\ifdefined\hypertarget\hypertarget{sourcepointer-1}{}\fi
\noindent\textit{Source reference.} Original lemma 0FSG: sources/stacks-sdga/pinned-original-source.tex, source lines 3129--3137. Cited by the retained statement or proof but not part of the collected spans; supplied only as a navigation pointer. The argument does not depend on its content.\par\medskip
\ifdefined\hypertarget\hypertarget{sourcepointer-2}{}\fi
\noindent\textit{Source reference.} Original lemma 0FSC: sources/stacks-sdga/pinned-original-source.tex, source lines 2983--2991. Cited by the retained statement or proof but not part of the collected spans; supplied only as a navigation pointer. The argument does not depend on its content.\par\medskip
\ifdefined\hypertarget\hypertarget{sourcepointer-3}{}\fi
\noindent\textit{Source reference.} Original lemma 0FSD: sources/stacks-sdga/pinned-original-source.tex, source lines 3022--3029. Cited by the retained statement or proof but not part of the collected spans; supplied only as a navigation pointer. The argument does not depend on its content.\par\medskip
\ifdefined\hypertarget\hypertarget{sourcepointer-4}{}\fi
\noindent\textit{Source reference.} Original lemma 0FSF: sources/stacks-sdga/pinned-original-source.tex, source lines 3098--3126. Cited by the retained statement or proof but not part of the collected spans; supplied only as a navigation pointer. The argument does not depend on its content.\par\medskip
\begin{lemma}
\label{lemma-resolve}
Let $(\mathcal{C}, \mathcal{O})$ be a ringed site.
Let $\mathcal{A}$ be a sheaf of differential graded algebras
on $(\mathcal{C}, \mathcal{O})$. Let $\mathcal{M}$
be a differential graded $\mathcal{A}$-module. There exists a homomorphism
$\mathcal{P} \to \mathcal{M}$ of differential graded $\mathcal{A}$-modules
with the following properties
\begin{enumerate}
\item $\mathcal{P} \to \mathcal{M}$ is a quasi-isomorphism, and
\item $\mathcal{P}$ is good.
\end{enumerate}
\end{lemma}

\begin{proof}[First proof]
Let $\mathcal{S}_0$ be the sheaf of graded sets
(Remark \ref{remark-sheaf-graded-sets})
whose degree $n$ part is $\Ker(\text{d}_\mathcal{M}^n)$.
Consider the homomorphism of differential graded modules
$$
\mathcal{P}_0 = \mathcal{A}[\mathcal{S}_0] \longrightarrow \mathcal{M}
$$
where the left hand side is as in Remark \ref{remark-sheaf-graded-sets}
and the map sends a local section $s$ of $\mathcal{S}_0$
to the corresponding local section of $\mathcal{M}^{\deg(s)}$
(which is in the kernel of the differential, so our map is
a map of differential graded modules indeed). By construction the
induced maps on cohomology sheaves $H^n(\mathcal{P}_0) \to H^n(\mathcal{M})$
are surjective. We are going to inductively construct maps
$$
\mathcal{P}_0 \to \mathcal{P}_1 \to \mathcal{P}_2 \to \ldots \to \mathcal{M}
$$
Observe that of course $H^*(\mathcal{P}_i) \to H^*(\mathcal{M})$
will be surjective for all $i$.
Given $\mathcal{P}_i \to \mathcal{M}$ denote $\mathcal{S}_{i + 1}$
the sheaf of graded sets whose degree $n$ part is
$$
\Ker(\text{d}_{\mathcal{P}_i}^{n + 1})
\times_{\mathcal{M}^{n + 1}, \text{d}}
\mathcal{M}^n
$$
Then we set
$$
\mathcal{P}_{i + 1} = \mathcal{P}_i \oplus \mathcal{A}[\mathcal{S}_{i + 1}]
$$
as graded $\mathcal{A}$-module with differential and map to
$\mathcal{M}$ defined as follows
\begin{enumerate}
\item for local sections of $\mathcal{P}_i$ use the differential on
$\mathcal{P}_i$ and the given map to $\mathcal{M}$,
\item for a local section $s = (p, m)$ of $\mathcal{S}_{i + 1}$ we set
$\text{d}(s)$ equal to $p$ viewed as a section of $\mathcal{P}_i$
of degree $\deg(s) + 1$ and we map $s$ to $m$ in $\mathcal{M}$, and
\item extend the differential uniquely so that the Leibniz rule holds.
\end{enumerate}
This makes sense because $\text{d}(m)$ is the image of $p$ and
$\text{d}(p) = 0$.
Finally, we set $\mathcal{P} = \colim \mathcal{P}_i$ with the
induced map to $\mathcal{M}$.

\medskip\noindent
The map $\mathcal{P} \to \mathcal{M}$ is a quasi-isomorphism:
we have $H^n(\mathcal{P}) = \colim H^n(\mathcal{P}_i)$
and for each $i$ the map $H^n(\mathcal{P}_i) \to H^n(\mathcal{M})$
is surjective with kernel annihilated  by the map
$H^n(\mathcal{P}_i) \to H^n(\mathcal{P}_{i + 1})$ by construction.
Each $\mathcal{P}_i$ is good because $\mathcal{P}_0$ is good
by Lemma \ref{lemma-free-graded-module-good} and each $\mathcal{P}_{i + 1}$
is in the middle of the admissible short exact sequence
$0 \to \mathcal{P}_i \to \mathcal{P}_{i + 1} \to
\mathcal{A}[\mathcal{S}_{i + 1}] \to 0$
whose outer terms are good by induction. Hence
$\mathcal{P}_{i + 1}$ is good by
Lemma \ref{lemma-good-admissible-ses}.
Finally, we conclude that $\mathcal{P}$ is good by
Lemma \ref{lemma-good-direct-sum}.
\end{proof}

\end{document}
